\chapter{Системы отсчёта, свет и геометрия пространства-времени}

\section{Курс вновь начинается с относительности}

В этой четверти мы возвращаемся к классической физике, но не к ньютоновской картине, в которой пространство и время независимы и универсальны. Непосредственными темами будут специальная теория относительности и классическая теория поля, особенно электромагнетизм. Квантовая механика использоваться не будет. Предыдущий курс механики дал принцип действия, а также гамильтоново и лагранжево описания; курс по квантовой механике дал её основные принципы, хотя даже гармонический осциллятор пришлось отложить до следующего возвращения к нему. Свет тоже почти не обсуждался. Теперь этот пробел восполняется, потому что свет по самой своей природе релятивистичен.

Большинство в аудитории уже сталкивалось со специальной теорией относительности. Тем не менее поднятие рук показывает, что полезно заново выстроить её с самого начала. Это и есть правильный выбор. Преобразование Лоренца не должно появляться как формула, которую нужно заучить. Оно должно возникнуть потому, что операциональные измерения пространства и времени не могут одновременно удовлетворять и галилеевой кинематике, и наблюдаемому закону распространения света.

\subsection{Координаты как физическая аппаратура}

Начнём с декартовой системы: начала координат, трёх ориентированных осей и координат \(x,y,z\). В конкретной реализации пространство было бы заполнено жёсткой решёткой из метровых стержней. Чтобы задавать событие, а не просто место, каждая точка решётки снабжается ещё и часами. Тогда событие обозначается
\begin{equation}
(x,y,z,t).
\end{equation}
Часы должны быть синхронизированы некоторой явно указанной процедурой; именно это требование вскоре окажется тонким местом конструкции.

\begin{definition}
Система отсчёта — это пространственная координатная решётка вместе с распределённой системой синхронизированных часов. Инерциальная система отсчёта — это невращающаяся, неускоренная система, в которой всякая свободная частица движется с постоянной скоростью.
\end{definition}

Две системы, движущиеся равномерно друг относительно друга, приписывают одному и тому же событию разные координаты. Точка, покоящаяся в движущейся системе, имеет постоянную координату \(x'\), хотя её координата без штриха меняется. В ньютоновской физике пространственное соотношение для движения вдоль направления \(x\) имеет вид
\begin{equation}
x'=x-vt.
\label{eq:sr1-galilean-space}
\end{equation}
Менее заметное ньютоновское предположение состоит в том, что все правильно выставленные часы остаются универсально синхронизированными:
\begin{equation}
t'=t.
\label{eq:sr1-absolute-time}
\end{equation}

\subsection{Принцип относительности}

Относительность началась не с Эйнштейна. Галилей и Ньютон уже пользовались принципом, согласно которому законы физики имеют один и тот же вид во всякой инерциальной системе отсчёта. Закон Ньютона \(F=ma\), закон тяготения обратных квадратов и сила Кулона не приобретают новой формы только потому, что лаборатория движется равномерно.

Обычная проверка — закрытый железнодорожный вагон. Если поезд идёт плавно и из окон ничего не видно снаружи, жонглирование происходит точно так же, как в неподвижной комнате. Шары поднимаются, падают и обычным образом возвращаются в руку. Эксперимент, целиком заключённый в вагоне, не может обнаружить его равномерную скорость. Ускорение или вращение можно обнаружить внутренними средствами; равномерное поступательное движение — нельзя.

Эйнштейн добавляет к этому более старому принципу один физический закон: свет в вакууме распространяется со скоростью
\begin{equation}
c \simeq 1.86\times 10^5\ {\rm miles/s}
   \simeq 3.00\times 10^8\ {\rm m/s}.
\end{equation}
Если это закон природы и законы одинаковы во всякой инерциальной системе отсчёта, то каждый инерциальный наблюдатель должен измерять одно и то же \(c\).

\section{Почему преобразование Галилея не работает}

Одного пространственного измерения достаточно, чтобы выявить противоречие. Нарисуем диаграмму пространства-времени с горизонтальной координатой \(x\) и вертикальной координатой \(t\). Световой импульс, испущенный из начала координат и движущийся вправо, описывается уравнением
\begin{equation}
x=ct.
\label{eq:sr1-light-right}
\end{equation}
Система, движущаяся вправо со скоростью \(v\), имеет своё начало координат на
\begin{equation}
x=vt,
\qquad
x'=0.
\label{eq:sr1-moving-origin}
\end{equation}

\begin{figure}[htbp]
\centering
\includegraphics[width=0.90\textwidth]{lecture_01_figure_01.png}
\caption{Идущий вправо световой луч и движущееся начало координат \(x=vt\), которое наблюдатель в штрихованной системе обозначает \(x'=0\).}
\lectureframe{00:15:24}
\end{figure}

Применим галилеевы правила \eqref{eq:sr1-galilean-space}--\eqref{eq:sr1-absolute-time} к световому лучу:
\begin{align}
x'&=x-vt=(c-v)t,\\
t'&=t.
\end{align}
Следовательно, наблюдатель в штрихованной системе заключил бы, что
\begin{equation}
x'=(c-v)t',
\end{equation}
так что свет, идущий вправо, имел бы скорость \(c-v\). Для луча, идущего влево, \(x=-ct\),
\begin{equation}
x'=-(c+v)t',
\end{equation}
скорость которого по модулю равна \(c+v\).

\begin{figure}[htbp]
\centering
\includegraphics[width=0.90\textwidth]{lecture_01_figure_02.png}
\caption{Галилеева подстановка предсказывает \(c-v\) для светового луча, идущего вправо, что противоречит не зависящей от системы отсчёта скорости света.}
\lectureframe{00:17:23}
\end{figure}

Обычное сложение скоростей — превосходное приближение для медленного механического движения, но измерения света не дают предсказанных \(c-v\) и \(c+v\). Пересмотреть нужно закон преобразования, а не постулат о свете. Главным подозреваемым является кажущееся безобидным утверждение \(t'=t\): оно утверждает абсолютное понятие одновременности.

\section{Синхронизация удалённых часов}

Чтобы сказать, что двое часов в разных местах показывают одно и то же время, требуется эксперимент. Поместим третьего наблюдателя в их середину, установив равенство расстояний метровыми стержнями. Пусть часы на обоих концах испускают вспышки, когда каждые из них показывают полдень. Если вспышки достигают середины одновременно, будем считать часы синхронизированными.

Это не философское определение, наложенное на физику извне. Это воспроизводимая операция, составленная из часов, расстояний и световых сигналов.

Теперь пусть некоторый наблюдатель движется мимо середины. Вспышки с концов, которые приходят одновременно в неподвижную середину, не приходят одновременно к движущемуся наблюдателю: пока свет находится в пути, он движется навстречу одной вспышке и удаляется от другой. Поэтому события, которые неподвижные часы считают одновременными, не являются одновременными согласно столь же корректной процедуре синхронизации в движущейся системе.

В нештрихованной системе линия одновременных событий задаётся как
\begin{equation}
t=\text{constant}.
\end{equation}
Соответствующая линия в штрихованной системе будет наклонена на диаграмме \(x\)-\(t\). Прежде чем спрашивать формулы преобразования, можно построить этот наклон с помощью световых сигналов.

\section{Построение движущейся системы с помощью света}

Положим \(c=1\). Можно измерять расстояние в световых секундах, если время измеряется в секундах, или измерять время по времени прохождения света через одну выбранную единицу длины. При любом из этих соглашений световые лучи имеют вид
\begin{equation}
x=\pm t,
\label{eq:sr1-light-units}
\end{equation}
и, следовательно, дают линии под углом \(45^\circ\) на одинаково масштабированных осях.

Сначала проведём световые лучи. Затем проведём семейство наблюдателей, покоящихся в движущейся системе:
\begin{equation}
x=vt,\qquad x=vt+1,\qquad x=vt+2.
\label{eq:sr1-moving-family}
\end{equation}
Их мировые линии параллельны, потому что наблюдатели сохраняют фиксированные расстояния в своей собственной инерциальной решётке.

\begin{figure}[htbp]
\centering
\includegraphics[width=0.90\textwidth]{lecture_01_figure_03.png}
\caption{Световые лучи под углом \(45^\circ\) и параллельные мировые линии \(x=vt\), \(x=vt+1\) и \(x=vt+2\) движущейся координатной решётки.}
\lectureframe{00:29:55}
\end{figure}

\subsection{Фред, Мэри и Сеймур}

Пусть Фред находится на \(x=vt\), Мэри — на \(x=vt+1\), а Сеймур — на \(x=vt+2\). В начале координат \(O\) Фред посылает направленный вправо световой сигнал к Мэри. Сеймур посылает сигнал влево так, чтобы обе вспышки достигли Мэри в одном и том же событии \(a\). Согласно правилу синхронизации по средней точке в движущейся системе, событие испускания Сеймура \(b\) одновременно с испусканием Фреда в \(O\). Поэтому прямая через \(O\) и \(b\) есть \(t'=0\).

Событие \(a\) лежит и на
\begin{equation}
x=t
\end{equation}
и на
\begin{equation}
x=vt+1.
\end{equation}
Следовательно,
\begin{equation}
t_a=x_a=\frac{1}{1-v}.
\label{eq:sr1-event-a}
\end{equation}

Идущий влево световой луч, проходящий через \(a\), удовлетворяет
\begin{equation}
x+t=x_a+t_a=\frac{2}{1-v}.
\label{eq:sr1-ab-light}
\end{equation}
Событие \(b\) также лежит на мировой линии Сеймура,
\begin{equation}
x-vt=2.
\label{eq:sr1-seymour-line}
\end{equation}
Решая \eqref{eq:sr1-ab-light} и \eqref{eq:sr1-seymour-line}, получаем
\begin{align}
t_b&=\frac{2v}{1-v^2},\\
x_b&=\frac{2}{1-v^2}.
\label{eq:sr1-event-b}
\end{align}
Поэтому наклон прямой, соединяющей \(O\) и \(b\), равен
\begin{equation}
\frac{t_b}{x_b}=v,
\end{equation}
так что линия одновременности движущейся системы имеет вид
\begin{equation}
t=vx,
\qquad
t'=0.
\label{eq:sr1-moving-simultaneity}
\end{equation}

\begin{figure}[htbp]
\centering
\includegraphics[width=0.93\textwidth]{lecture_01_figure_04.png}
\caption{Построение синхронизации Фред--Мэри--Сеймур. Световые сигналы, приходящие одновременно, определяют \(O\) и \(b\) как одновременные в движущейся системе, что даёт \(t=vx\).}
\lectureframe{00:44:38}
\end{figure}

Две оси штрихованной системы теперь записываются параллельным образом:
\begin{align}
x'=0 &\quad\Longleftrightarrow\quad x=vt,\\
t'=0 &\quad\Longleftrightarrow\quad t=vx.
\end{align}
Они являются зеркальными отражениями друг друга относительно идущей вправо световой линии \(x=t\). Геометрия уже заменила универсальную одновременность.

\begin{classroomqa}
\audiencequestion{Что происходит с этим построением, когда \(v=1\) или \(v=-1\)?}
\lecturerresponse{Знаменатели в \eqref{eq:sr1-event-a}--\eqref{eq:sr1-event-b} обращаются в нуль, и соответствующие точки пересечения уходят на бесконечность. При \(\lvert v\rvert=1\) материальный наблюдатель двигался бы по световому лучу и не смог бы выполнить обычную процедуру обмена через середину. Это сингулярное поведение — первое геометрическое предупреждение о том, что \(c\) является предельной скоростью. \lecturetimestamp{00:40:43}}
\end{classroomqa}

\section{Вывод преобразования Лоренца}

Требование, чтобы \(x'=0\) при \(x=vt\), означает, что в \(x'\) должен входить множитель \(x-vt\). Аналогично, требование, чтобы \(t'=0\) при \(t=vx\), означает наличие множителя \(t-vx\). Однородность пространства-времени делает преобразование линейным, поэтому запишем
\begin{align}
x'&=(x-vt)f(v),\\
t'&=(t-vx)g(v).
\label{eq:sr1-linear-ansatz}
\end{align}

\begin{figure}[htbp]
\centering
\includegraphics[width=0.90\textwidth]{lecture_01_figure_05.png}
\caption{Линейное преобразование строится из двух условий обращения в нуль, с зависящими от скорости масштабными множителями \(f(v)\) и \(g(v)\).}
\lectureframe{00:49:35}
\end{figure}

Для светового луча, идущего вправо, \(x=t\), инвариантность скорости света требует \(x'=t'\). Тогда уравнение~\eqref{eq:sr1-linear-ansatz} принимает вид
\begin{align}
x'&=(1-v)x f(v),\\
t'&=(1-v)x g(v),
\end{align}
и, следовательно,
\begin{equation}
f(v)=g(v).
\label{eq:sr1-f-equals-g}
\end{equation}
Таким образом, остаётся один общий масштабный множитель:
\begin{align}
x'&=(x-vt)f(v),\\
t'&=(t-vx)f(v).
\label{eq:sr1-direct-unscaled}
\end{align}
По существу, это тот же экономный ход рассуждений, которым Эйнштейн воспользовался в своей первой статье по теории относительности: операционное определение одновременности и инвариантность скорости света определяют преобразование.

\subsection{Взаимность фиксирует масштабный множитель}

Нет никакой выделенной \emph{действительно движущейся} инерциальной системы отсчёта. Если штрихованная система движется со скоростью \(v\) относительно нештрихованной, то нештрихованная движется со скоростью \(-v\) относительно штрихованной. Поэтому обратное отображение должно иметь тот же вид, но с обратным знаком у \(v\):
\begin{align}
x&=(x'+vt')f(v),\\
t&=(t'+vx')f(v).
\label{eq:sr1-inverse-unscaled}
\end{align}

\begin{figure}[htbp]
\centering
\includegraphics[width=0.90\textwidth]{lecture_01_figure_06.png}
\caption{Из взаимности следует обратное преобразование: знак относительной скорости меняется, а тот же масштабный множитель сохраняется.}
\lectureframe{00:55:12}
\end{figure}

Подставим \eqref{eq:sr1-direct-unscaled} в первое из уравнений \eqref{eq:sr1-inverse-unscaled}:
\begin{align}
x
&=\left[(x-vt)f(v)+v(t-vx)f(v)\right]f(v)\\
&=x(1-v^2)f(v)^2.
\end{align}
Поскольку эта композиция должна быть тождественным преобразованием,
\begin{equation}
(1-v^2)f(v)^2=1,
\end{equation}
и, следовательно,
\begin{equation}
f(v)=\pm\frac{1}{\sqrt{1-v^2}}.
\label{eq:sr1-factor-roots}
\end{equation}

\begin{figure}[htbp]
\centering
\includegraphics[width=0.90\textwidth]{lecture_01_figure_07.png}
\caption{Композиция прямого и обратного отображений фиксирует величину общего множителя \(f(v)\).}
\lectureframe{00:59:31}
\end{figure}

\begin{classroomqa}
\audiencequestion{Почему в \eqref{eq:sr1-factor-roots} отбрасывают отрицательный квадратный корень?}
\lecturerresponse{При нулевой относительной скорости преобразование должно непрерывно переходить в тождественное преобразование, \(x'=x\) и \(t'=t\). Отрицательная ветвь, напротив, даёт \(x'=-x\) и \(t'=-t\), то есть обращение обеих ориентаций координат. Её выбор вводил бы скачкообразное отражение вместо непрерывно движущейся системы отсчёта. \lecturetimestamp{00:59:27}}
\end{classroomqa}

При выборе положительной ветви преобразование Лоренца в единицах \(c=1\) имеет вид
\begin{align}
x'&=\frac{x-vt}{\sqrt{1-v^2}},&
t'&=\frac{t-vx}{\sqrt{1-v^2}}.
\label{eq:sr1-lorentz-c1}
\end{align}

\begin{figure}[htbp]
\centering
\includegraphics[width=0.90\textwidth]{lecture_01_figure_08.png}
\caption{Полное преобразование Лоренца в единицах, в которых скорость света равна единице.}
\lectureframe{01:02:05}
\end{figure}

\subsection{Восстановление обычных единиц}

Обычные единицы восстанавливаются с помощью размерностного анализа. Знаменатель может зависеть только от безразмерного отношения \(v^2/c^2\), тогда как \(vx\) в уравнении для времени должен превратиться в \(vx/c^2\):
\begin{align}
x'&=\gamma(x-vt),&
t'&=\gamma\left(t-\frac{vx}{c^2}\right),
\label{eq:sr1-lorentz-dimensional}\\
\gamma&=\frac{1}{\sqrt{1-v^2/c^2}}.
\label{eq:sr1-gamma}
\end{align}
Обратное преобразование имеет вид
\begin{align}
x&=\gamma(x'+vt'),&
t&=\gamma\left(t'+\frac{vx'}{c^2}\right).
\label{eq:sr1-lorentz-inverse-dimensional}
\end{align}

\begin{figure}[htbp]
\centering
\includegraphics[width=0.90\textwidth]{lecture_01_figure_09.png}
\caption{Формулы Лоренца с восстановленным \(c\): \(v^2/c^2\) задаёт масштаб, а \(vx/c^2\) вносит координату положения в выражение для времени.}
\lectureframe{01:05:01}
\end{figure}

При \(v\ll c\),
\begin{equation}
\gamma=1+\frac{1}{2}\frac{v^2}{c^2}+O\!\left(\frac{v^4}{c^4}\right),
\end{equation}
и преобразование сводится к виду
\begin{equation}
x'\approx x-vt,
\qquad
t'\approx t.
\end{equation}
При земных скоростях эти отношения обычно очень малы, что объясняет успех ньютоновской кинематики. Вблизи \(c\) поправки становятся большими.

Движение было выбрано вдоль оси \(x\). Поперечное измерительное плечо не может отличить одну инерциальную систему отсчёта от другой, не нарушая взаимности, поэтому
\begin{equation}
y'=y,
\qquad
z'=z.
\end{equation}

\section{Измерение движущегося стержня}

Под длиной понимается расстояние между двумя концами, измеренное в один и тот же момент времени в системе отсчёта, проводящей измерение. Это не означает кажущийся угловой размер сфотографированного объекта, поскольку задержки распространения света создают отдельную оптическую задачу.

Пусть стержень покоится в нештрихованной системе отсчёта и имеет концы
\begin{equation}
x=0,
\qquad
x=1.
\end{equation}
Наблюдатель в штрихованной системе измеряет события на концах на одном срезе одновременности по штрихованному времени, скажем при \(t'=0\). Для дальнего конца \(x=1\),
\begin{equation}
t'=0
\quad\Longrightarrow\quad
t=\frac{vx}{c^2}.
\end{equation}
Временно возвращаясь к \(c=1\), для этого события \(t=v\), и
\begin{equation}
x'=\frac{x-vt}{\sqrt{1-v^2}}
   =\frac{1-v^2}{\sqrt{1-v^2}}
   =\sqrt{1-v^2}.
\end{equation}

\begin{figure}[htbp]
\centering
\includegraphics[width=0.90\textwidth]{lecture_01_figure_10.png}
\caption{Мировые линии \(x=0\) и \(x=1\) стержня, покоящегося в нештрихованной системе отсчёта; движущийся наблюдатель должен сравнивать их на одном срезе \(t'\).}
\lectureframe{01:11:14}
\end{figure}

\begin{figure}[htbp]
\centering
\includegraphics[width=0.90\textwidth]{lecture_01_figure_11.png}
\caption{Одновременное в штрихованной системе измерение концов даёт коэффициент сокращения \(\sqrt{1-v^2}\) в единицах \(c=1\).}
\lectureframe{01:16:03}
\end{figure}

После восстановления единиц и для произвольной собственной длины \(L_0\),
\begin{equation}
L=\frac{L_0}{\gamma}
 =L_0\sqrt{1-\frac{v^2}{c^2}}.
\label{eq:sr1-length-contraction}
\end{equation}

\begin{classroomqa}
\audiencequestion{Если дальний конец стержня находится в \(x=1\), не должна ли его координата также быть \(x'=1\)?}
\lecturerresponse{Нет. Стержень покоится в нештрихованной системе отсчёта, поэтому \(x=0\) и \(x=1\) обозначают там его фиксированные концы. Штрихованная координата дальнего конца меняется со временем. Её значение на срезе измерения \(t'=0\) как раз и вычисляется преобразованием Лоренца. \lecturetimestamp{01:14:24}}
\end{classroomqa}

\begin{classroomqa}
\audiencequestion{Как может каждый инерциальный наблюдатель последовательно утверждать, что стержень другого наблюдателя короче?}
\lecturerresponse{Они измеряют разные пары событий на концах. Нештрихованному наблюдателю требуется равенство \(t\); штрихованному наблюдателю требуется равенство \(t'\). Поскольку одновременность зависит от системы отсчёта, эти две пары концовых событий не являются одной и той же парой событий. Поэтому взаимность не приводит к противоречию. \lecturetimestamp{01:16:28}}
\end{classroomqa}

Операционную версию стоит сформулировать явно. Расставьте вдоль измеряющей системы отсчёта синхронизированные часы и помощников. В заранее заданный местный момент времени каждый помощник записывает положение конца, проходящего рядом. Разность этих двух положений, зафиксированных одновременно, даёт длину стержня в этой системе. Ничто не зависит от того, как стержень при этом выглядит.

\begin{classroomqa}
\audiencequestion{Когда \(c=1\), обязан ли стержень единичной длины иметь длину в одну световую секунду?}
\lecturerresponse{Только если время измеряется в секундах. Вместо этого можно использовать метры для длины и измерять время временем, за которое свет проходит один метр. Приравнивание \(c=1\) связывает единицы времени и длины; оно не предписывает, какую физическую единицу нужно выбрать первой. \lecturetimestamp{01:21:17}}
\end{classroomqa}

\section{Движущиеся часы и вопрос о близнецах}

Теперь проследим за часами, покоящимися в штрихованной системе. Вдоль их мировой линии,
\begin{equation}
x'=0.
\end{equation}
Сравним с началом координат событие, в котором они показывают \(t'=1\). Обратное преобразование даёт
\begin{equation}
t=\gamma\left(t'+\frac{vx'}{c^2}\right)=\gamma.
\end{equation}
Следовательно, одной единице прошедшего времени на движущихся часах соответствует \(\gamma\) единиц нештрихованного координатного времени:
\begin{equation}
\Delta t=\gamma\,\Delta\tau,
\qquad
\Delta\tau=\frac{\Delta t}{\gamma}.
\label{eq:sr1-time-dilation}
\end{equation}

\begin{figure}[htbp]
\centering
\includegraphics[width=0.90\textwidth]{lecture_01_figure_12.png}
\caption{Часы, переносимые вдоль \(x'=0\), показывают меньше прошедшего времени между теми же двумя событиями, чем координатное время покоящейся системы: \(t=\gamma t'\).}
\lectureframe{01:26:42}
\end{figure}

Этот эффект относится ко всяким физическим часам, а не только к механизму, который называется часами. Наручные часы, атомный переход, метаболизм и биологическое старение — всё подчиняется одному и тому же собственному времени. Дешёвые часы, купленные на улице в Нью-Йорке, не избегут релятивистских эффектов из-за того, что это плохой Rolex; они могут идти неточно, но какой бы процесс в них ни использовался, он всё равно подчиняется той же кинематике.

\begin{classroomqa}
\audiencequestion{Если каждый инерциальный наблюдатель видит, что часы другого идут медленно, почему путешествующий близнец возвращается моложе?}
\lecturerresponse{Путешествие туда и обратно не является сравнением двух глобально эквивалентных инерциальных мировых линий. Путешественник должен развернуться и перейти в другую инерциальную систему отсчёта, тогда как оставшийся дома близнец остаётся на одной инерциальной мировой линии. Разворот нарушает симметрию обмена, поэтому накопленные собственные времена могут различаться. \lecturetimestamp{01:26:37}}
\end{classroomqa}

\begin{classroomqa}
\audiencequestion{Является ли именно ускорение во время разворота причиной разницы в возрасте?}
\lecturerresponse{Ускорение отмечает место, где нарушается симметрия, но разность возрастов вычисляется интегрированием собственного времени вдоль полных траекторий. Разворот можно сделать кратким или рассматривать как переход между инерциальными участками; неравные пути в пространстве-времени всё равно имеют неравные собственные времена. \lecturetimestamp{01:29:20}}
\end{classroomqa}

\section{Инвариантная геометрия}

Диаграмма замедления времени выглядит парадоксально, если читать её в евклидовой геометрии. Наклонный путь кажется гипотенузой, однако ему соответствует меньше собственного времени, чем вертикальному пути. Разгадка не сводится к алгебраическому фокусу. У пространства-времени другая метрика.

На евклидовой плоскости поворот изменяет \(x\) и \(y\), но сохраняет
\begin{equation}
x'^2+y'^2=x^2+y^2.
\label{eq:sr1-euclidean-invariant}
\end{equation}
Естественно проверить, сохраняет ли преобразование Лоренца \(t^2+x^2\). Нет. В единицах \(c=1\),
\begin{align}
t'^2
&=\frac{t^2-2vtx+v^2x^2}{1-v^2},\\
x'^2
&=\frac{x^2-2vtx+v^2t^2}{1-v^2}.
\end{align}
При сложении получаются неверные коэффициенты. Если же вместо этого вычесть, смешанные члены сократятся, и получится
\begin{align}
t'^2-x'^2
&=
\frac{t^2-2vtx+v^2x^2-x^2+2vtx-v^2t^2}{1-v^2}\\
&=t^2-x^2.
\label{eq:sr1-interval-invariance}
\end{align}

\begin{figure}[htbp]
\centering
\includegraphics[width=0.90\textwidth]{lecture_01_figure_13.png}
\caption{Вычитание квадратов преобразованных координат сокращает смешанные члены и выявляет инвариант \(t^2-x^2\).}
\lectureframe{01:38:07}
\end{figure}

Координаты — это зависящие от системы отсчёта компоненты. Комбинация
\begin{equation}
\Delta s^2=c^2\Delta t^2-\Delta x^2
\label{eq:sr1-spacetime-interval}
\end{equation}
одинакова во всех инерциальных системах отсчёта. При большем числе пространственных измерений,
\begin{equation}
\Delta s^2
=c^2\Delta t^2-\Delta x^2-\Delta y^2-\Delta z^2.
\end{equation}

\begin{classroomqa}
\audiencequestion{Почему правомерно отказаться от обычного геометрического правила, согласно которому гипотенуза длиннее?}
\lecturerresponse{Это правило относится к положительно определённой евклидовой метрике. Здесь координаты были операционально определены с помощью линеек, часов и световых сигналов, а их преобразование сохраняет разность квадратов. Геометрия Минковского не постулируется, чтобы избежать парадокса; она выводится из этих измерений.\lecturetimestamp{01:43:21}}
\end{classroomqa}

\subsection{Собственное время}

Для двух событий на мировой линии часов выберем систему покоя этих часов. В ней \(\Delta x'=0\), поэтому из инвариантности следует
\begin{equation}
c^2\Delta\tau^2
=c^2\Delta t^2-\Delta x^2,
\end{equation}
или
\begin{equation}
\Delta\tau
=\sqrt{\Delta t^2-\frac{\Delta x^2}{c^2}}.
\label{eq:sr1-proper-time}
\end{equation}
Это и есть время, которое часы действительно отсчитывают между двумя событиями.

\begin{figure}[htbp]
\centering
\includegraphics[width=0.90\textwidth]{lecture_01_figure_14.png}
\caption{В системе покоя часов пространственное разделение равно нулю, поэтому инвариантный интервал становится квадратом их собственного времени.}
\lectureframe{01:41:28}
\end{figure}

\begin{classroomqa}
\audiencequestion{Должно ли собственное время быть \(t^2-x^2/c^2\) или его квадратным корнем?}
\lecturerresponse{Разность квадратов — это квадрат собственного времени. Прошедшее по часам время равно \(\tau=\sqrt{t^2-x^2/c^2}\) для разделений, измеряемых от общего начала. Пропущенный квадратный корень в первоначальной формулировке необходимо восстановить.\lecturetimestamp{01:51:02}}
\end{classroomqa}

\subsection{Времениподобное, нулевое и пространственноподобное разделение}

Два различных события могут иметь нулевой интервал. В единицах, где \(c=1\),
\begin{equation}
\Delta t^2-\Delta x^2=0
\quad\Longleftrightarrow\quad
\Delta x=\pm\Delta t.
\end{equation}
Такие события соединяет световой луч.

\begin{figure}[htbp]
\centering
\includegraphics[width=0.90\textwidth]{lecture_01_figure_15.png}
\caption{События с нулевым разделением лежат на \(x=\pm t\), где пространственно-временной интервал и собственное время светоподобного пути равны нулю.}
\lectureframe{01:46:40}
\end{figure}

Если бы свет мог нести обычные часы, между испусканием и поглощением не произошло бы ни одного тика собственного времени. Речь идёт не о допустимой системе покоя для света; это предельное утверждение о том, что любой нулевой отрезок имеет нулевое собственное время.

Световые лучи делят пространство-время на три инвариантных класса:
\begin{align}
c^2\Delta t^2-\Delta x^2&>0
&&\text{timelike},\\
c^2\Delta t^2-\Delta x^2&=0
&&\text{null},\\
c^2\Delta t^2-\Delta x^2&<0
&&\text{spacelike}.
\end{align}
Реальные часы могут соединить времениподобно разделённые события, и положительный интервал определяет собственное время этих часов. Ни один досветовой объект не может соединить пространственноподобно разделённые события. Для пространственноподобной пары можно выбрать систему отсчёта, в которой события одновременны, и определить собственное расстояние
\begin{equation}
\ell_{\rm proper}
=\sqrt{\Delta x^2-c^2\Delta t^2}.
\end{equation}

\begin{figure}[htbp]
\centering
\includegraphics[width=0.90\textwidth]{lecture_01_figure_16.png}
\caption{Световой конус отделяет времениподобные пары, которые часы могут соединить, от пространственноподобных пар, которые не может соединить ни один досветовой объект.}
\lectureframe{01:48:47}
\end{figure}

\begin{classroomqa}
\audiencequestion{Когда инвариант следует называть собственным временем, а когда — собственной длиной?}
\lecturerresponse{Для времениподобной пары существует система покоя, в которой оба события происходят в одном месте, и интервал даёт собственное время. Для пространственноподобной пары существует система отсчёта, в которой оба происходят в одно и то же время, и инвариант противоположного знака даёт их собственное расстояние. Нулевое разделение является границей между ними.\lecturetimestamp{01:47:22}}
\end{classroomqa}

\begin{classroomqa}
\audiencequestion{Где должно стоять \(c\), если не используются единицы, в которых \(c=1\)?}
\lecturerresponse{Инвариант имеет вид \(c^2\Delta t^2-\Delta x^2\) и размерность квадрата длины. Деление на \(c^2\) даёт \(\Delta\tau^2=\Delta t^2-\Delta x^2/c^2\) с размерностью квадрата времени.\lecturetimestamp{01:49:56}}
\end{classroomqa}

\section{Симметрия Лоренца и её исторический смысл}

\begin{classroomqa}
\audiencequestion{Можно ли рассматривать преобразование Лоренца как вращение?}
\lecturerresponse{Да, но не как обычное евклидово вращение. Это гиперболическое вращение, сохраняющее \(t^2-x^2\). Исторически его можно описывать как вращение на мнимый угол; быстрота даёт современную вещественную параметризацию.\lecturetimestamp{01:51:40}}
\end{classroomqa}

Подход Эйнштейна был тесно связан с электромагнетизмом. Согласно его позднейшему воспоминанию, ещё подростком он задавался вопросом, что случится, если погнаться за световой волной. Галилеево вычитание скоростей заставило бы волну казаться неподвижной, но уравнения Максвелла не допускают соответствующей застывшей колебательной световой волны в вакууме. Этот конфликт указывал на кинематику, а не на какое-то особое механическое свойство измерительных стержней.

\begin{classroomqa}
\audiencequestion{Был ли аргумент Эйнштейна прежде всего ответом на эксперименты, подобные эксперименту Майкельсона--Морли?}
\lecturerresponse{Экспериментальный фон исторически имел значение, но сам Эйнштейн подчёркивал структуру уравнений Максвелла и мысленный эксперимент погони за световым лучом. Исторически спорно, насколько хорошо он знал конкретные результаты интерферометра или насколько на них опирался, поэтому ясный физический смысл состоит в конфликте симметрий между Максвеллом и преобразованиями Галилея.\lecturetimestamp{01:52:04}}
\end{classroomqa}

Уравнения Максвелла инвариантны относительно преобразований Лоренца, а не Галилея. Лоренц и Фицджеральд уже получили родственные формулы сокращения, сохраняя при этом эфир и интерпретируя сокращение как динамическую деформацию, вызванную движением сквозь него. В рамках такого взгляда этот эффект мог быть приближённым материальным откликом с дополнительными поправками.

Решающим шагом Эйнштейна было сделать принцип относительности и инвариантную скорость света точными кинематическими принципами. Тогда сокращение длины и замедление времени применимы к любому подходящему мерному стержню и любым подходящим часам, потому что они выражают геометрию, связывающую инерциальные системы отсчёта, а не особую силу, оказываемую эфиром. Следующая задача — сформулировать движение частицы в рамках этой геометрии.
